\section{积分的构造与极限定理}\label{sec:03-01}

\subsection{从简单函数到正积分：Daniell 视角}\label{subsec:3-1-1}
\index{简单函数}\index{非负积分}\index{Daniell--Young 积分}

测度已经回答了集合有多大，却还没有告诉我们函数图像有多少 ``质量''。在有限原子空间里，
答案没有悬念。若
\[
  \mu=\sum_{j=1}^N w_j\delta_{x_j},\qquad w_j\geq0,
\]
那么任何合理的积分都应复现加权和
\[
  \int f\,\dd\mu=\sum_{j=1}^Nw_jf(x_j).
  \tag{3.1.1}\label{eq:3-1-finite-observation}
\]
真正的问题是：当函数有无穷多个取值、空间也不再由有限个原子组成时，怎样从有限和走过去？

最先能够逐项计算的函数不是连续函数，而是只取有限多个值的函数。由集合到函数的扩张因而从这里开始。
这条路线与 Daniell--Young 的观点平行：先把 ``正线性'' 和 ``从下逼近'' 当成不可让步的规则，
再让极限定理说明所得对象确实稳定。这里从已有测度出发构造积分；从正泛函反向恢复测度则是
Daniell 表示定理所处理的逆问题。

\begin{definition}[非负简单函数的积分]\label{def:3-1-1-1}
设 $(X,\mathcal A,\mu)$ 为测度空间。沿用定义~\ref{def:2-2-1-simple-function}，把非负简单函数写成
\[
  s=\sum_{k=1}^m a_k\1_{E_k},
\]
其中 $a_k\in[0,\infty)$，$E_k\in\mathcal A$ 两两不交，并定义
\[
  \int_Xs\,\dd\mu:=\sum_{k=1}^ma_k\mu(E_k).
\]
这里以及下文均约定 $0\cdot\infty=0$。
\end{definition}

同一简单函数可以有许多写法。例如，值为零的集合可以任意拆分，两个取同一系数的集合也可以合并。
因此首先要确认上式算的是函数，而不是写法。

\begin{lemma}[简单函数积分表示无关]\label{lemma:3-1-1-1}
简单函数的积分不依赖于所取的不交值层表示。对非负简单函数 $s,t$ 与 $a,b\geq0$，还有
\[
  \int(as+bt)\,\dd\mu=a\int s\,\dd\mu+b\int t\,\dd\mu.
\]
\end{lemma}

\begin{proof}
设
\[
 s=\sum_{i=1}^m a_i\1_{E_i}=\sum_{j=1}^n b_j\1_{F_j}
\]
是两个两两不交表示。加入余集并令其系数为零后，可设两族各自分割 $X$。集合
$E_i\cap F_j$ 构成公共有限分割。若 $E_i\cap F_j\ne\varnothing$，在其中任取一点便有
$a_i=s(x)=b_j$；若交集为空，其测度为零。因此
\[
 \sum_i a_i\mu(E_i)
 =\sum_{i,j}a_i\mu(E_i\cap F_j)
 =\sum_{i,j}b_j\mu(E_i\cap F_j)
 =\sum_jb_j\mu(F_j).
\]
等式只涉及非负扩展实数，因而没有无穷减无穷的问题。对 $s,t$ 的分割取公共细分，
在每个小块上直接合并系数，即得所述正线性。
\end{proof}

\begin{example}[三个有限观测模型]\label{ex:3-1-1-1}
若 $\mu$ 是计数测度，则对非负简单函数 $s$，
\[
 \int s\,\dd\mu=\sum_{x\in X}s(x),
\]
其中右侧定义为所有有限子集和的上确界；由于 $s$ 只有有限多个正值层，这与上面的定义一致。
若 $\mu=\delta_x$，则 $\int s\,\dd\delta_x=s(x)$。更一般地，有限原子测度立即给出
\eqref{eq:3-1-finite-observation}。这些计算提示我们：对非负量，添加更多观测只能使答案增加，
所以下确界逼近不如从下逼近自然。
\end{example}

\begin{example}[Dirac 测度]\label{ex:3-1-1-2}
上例的单原子情形值得单独记住。对任意非负可测函数 $f$，稍后的定义将给出
\[
 \int f\,\dd\delta_x=f(x),
\]
包括 $f(x)=+\infty$ 的情形。换言之，Dirac 积分就是一次点值观测。
\end{example}

\begin{definition}[非负积分]\label{def:3-1-1-2}
对可测函数 $f:X\to[0,+\infty]$，定义
\[
 \int_X f\,\dd\mu
 :=\sup\left\{\int_Xs\,\dd\mu:
      s\text{ 为非负简单函数且 }0\leq s\leq f\right\}.
 \tag{3.1.2}\label{eq:3-1-positive-integral}
\]
对 $E\in\mathcal A$，记
\[
 \int_Ef\,\dd\mu:=\int_Xf\1_E\,\dd\mu.
\]
\end{definition}

定义中的集合非空，因为零函数总是候选者；其上确界允许为 $+\infty$。经验测度
$n^{-1}\sum_{j=1}^n\delta_{x_j}$ 下，$\int_Ef$ 最终将等于只对满足 $x_j\in E$ 的样本求平均，
所以集合下标只是对有限观测作筛选，并未偷偷引入另一种积分。

\begin{definition}[几乎处处与共同零集]\label{def:3-1-1-ae}
若存在可测零集 $N$，使性质 $P(x)$ 对每个 $x\in X\setminus N$ 成立，就称 $P$ 在
$\mu$-几乎处处意义下成立，简记为 $\mu$-a.e. 成立。数列 $f_n\to f$ a.e. 时，
要求有一个对所有 $n$ 共同有效的可测零集。可数多个零集之并仍为零集，正是数列可以共同化异常集的原因。
\end{definition}

定义~\ref{def:3-1-1-2} 看似要在一个巨大的函数族上取上确界。事实上，有一列完全明确的阶梯就够了。

\begin{lemma}[规范层级逼近]\label{lemma:3-1-1-layer}
对非负可测函数 $f$，令
\[
 s_n(f)=2^{-n}\bigl\lfloor 2^n(f\wedge n)\bigr\rfloor,
 \qquad n\geq1,
 \tag{3.1.3}\label{eq:3-1-canonical-simple}
\]
并约定 $\lfloor+\infty\rfloor=+\infty$，但由于先取了 $f\wedge n$，右侧始终有限。
则 $s_n(f)$ 是有限值非负简单函数，
\[
  s_n(f)\uparrow f,
  \qquad
  \int f\,\dd\mu=\sup_n\int s_n(f)\,\dd\mu.
  \tag{3.1.4}\label{eq:3-1-canonical-sup}
\]
此外，若简单函数 $t_n\uparrow t$，且 $t$ 也是有限值简单函数，则
$\int t_n\uparrow\int t$。
\end{lemma}

\begin{proof}
函数 $s_n(f)$ 的值属于有限集合
$\{0,2^{-n},\ldots,n\}$，其各个值层由 $f$ 的可测逆像构成，故它是简单函数。
对任意 $u\in[0,+\infty]$，先逐情形核对单调性。若 $u\le n$，则
\[
 \lfloor2^{n+1}u\rfloor=\lfloor2(2^nu)\rfloor
 \ge 2\lfloor2^nu\rfloor,
\]
从而
\[
 2^{-(n+1)}\lfloor2^{n+1}(u\wedge(n+1))\rfloor
 \ge 2^{-n}\lfloor2^n(u\wedge n)\rfloor.
\]
若 $u>n$，则
\[
 2^{-(n+1)}\lfloor2^{n+1}(u\wedge(n+1))\rfloor
 \ge 2^{-(n+1)}\lfloor2^{n+1}n\rfloor=n
 =2^{-n}\lfloor2^n(u\wedge n)\rfloor,
\]
所以同一不等式仍成立。再若 $u<\infty$，当 $n>u$ 时
\[
 0\le u-2^{-n}\lfloor2^nu\rfloor<2^{-n},
\]
故该数列趋于 $u$；若 $u=+\infty$，则其第 $n$ 项正好为 $n$。逐点应用便得
$s_n(f)\uparrow f$。

由定义，$\sup_n\int s_n(f)\leq\int f$。反过来，任取非负简单函数 $s\leq f$ 和
$0<c<1$。若 $a>0$ 遍历 $s$ 的有限多个正值，取 $n$ 足够大，使
$n\geq\max a$ 且 $2^{-n}\leq(1-c)\min a$。在 $s(x)=a$ 的值层上，
\[
 s_n(f)(x)\geq2^{-n}\lfloor2^na\rfloor\geq a-2^{-n}\geq ca,
\]
而在零值层上不等式也成立。因此 $cs\leq s_n(f)$，从而
\[
 c\int s\,\dd\mu\leq\sup_n\int s_n(f)\,\dd\mu.
\]
先令 $c\uparrow1$，再对 $s$ 取上确界，得到 \eqref{eq:3-1-canonical-sup}。
这里没有使用任何积分极限定理。

最后设 $t=\sum_{j=1}^r a_j\1_{E_j}$，其中 $a_j>0$ 且 $E_j$ 两两不交；零值层可略去。
固定 $0<c<1$，令 $E_{j,n}=E_j\cap\{t_n\geq ca_j\}$。由 $t_n\uparrow t$，
$E_{j,n}\uparrow E_j$。测度从下连续给出
\[
 \begin{aligned}
 \liminf_n\int t_n\,\dd\mu
 &\geq \lim_n\int c\sum_j a_j\1_{E_{j,n}}\,\dd\mu\\
 &=c\sum_j a_j\lim_n\mu(E_{j,n})
 =c\int t\,\dd\mu.
 \end{aligned}
\]
令 $c\uparrow1$；再结合 $t_n\leq t$ 所给的反向不等式，结论成立。
这一段只用了测度的从下连续性，而该性质已由可数可加性证明。
\end{proof}

\begin{example}[层级阶梯]\label{ex:3-1-1-3}
若 $f(x)=x$ 定义在 $[0,1]$ 上，则 $s_n(f)$ 在每个
$[k2^{-n},(k+1)2^{-n})$ 上取值 $k2^{-n}$。因此
\[
 \int_0^1s_n(f)\,\dd x
 =2^{-2n}\sum_{k=0}^{2^n-1}k
 =\frac12-2^{-n-1}.
\]
由 \eqref{eq:3-1-canonical-sup}，$\int_0^1x\,\dd x=1/2$。同一计算也说明，
定义中的上确界不是一件不可计算的装置：它把熟悉的下矩形和直接纳入其中。
\end{example}

\begin{definition}[正线性泛函]\label{def:3-1-1-3}
设 $\mathcal C$ 是一个非负函数锥。映射 $I:\mathcal C\to[0,+\infty]$ 若保持次序，且
\[
 I(af+bg)=aI(f)+bI(g)\qquad(a,b\geq0),
\]
则称为正线性泛函。扩展值等式仍采用 $0\cdot(+\infty)=0$，并且不作
$+\infty-(+\infty)$ 运算。
\end{definition}

下一命题说明，从下方简单函数取上确界确实保留了有限和的代数。正向不等式直接来自定义，
反向不等式则需要同时细分两个简单函数逼近。

\begin{proposition}[非负积分的代数性质]\label{proposition:3-1-1-2}
若 $f,g:X\to[0,+\infty]$ 可测且 $a,b\geq0$，则
\[
 \int(af+bg)\,\dd\mu
 =a\int f\,\dd\mu+b\int g\,\dd\mu.
 \tag{3.1.5}\label{eq:3-1-positive-linearity}
\]
若 $f\leq g$，则 $\int f\leq\int g$。
\end{proposition}

\begin{proof}
单调性直接来自 \eqref{eq:3-1-positive-integral}：$f$ 下方的每个简单函数也是 $g$ 下方的候选者。
若 $a>0$，映射 $s\mapsto as$ 把 $f$ 下方的简单函数与 $af$ 下方的简单函数一一对应，
所以 $\int af=a\int f$；$a=0$ 由约定成立。因而只须证明可加性。

若 $p\leq f$、$q\leq g$ 为非负简单函数，则 $p+q\leq f+g$，由引理
\ref{lemma:3-1-1-1}，
\[
 \int(f+g)\geq\int(p+q)=\int p+\int q.
\]
对非负扩展实数族有
$\sup_{p,q}(A_p+B_q)=\sup_pA_p+\sup_qB_q$：有限情形可分别逼近两个上确界，
任一上确界为无穷时则先把对应项取得任意大。因而分别对 $p,q$ 取上确界，得到
$\int(f+g)\geq\int f+\int g$。

为证另一方向，固定任意有限值简单函数 $s\leq f+g$，并令
\[
 u=s\wedge f,
 \qquad v=s-u.
\]
由于 $0\leq u\leq s$，$u,v$ 都是有限值可测函数，且 $u+v=s$、$u\leq f$。
若 $f\geq s$，则 $v=0$；若 $f<s$，则 $v=s-f\leq g$。故始终有 $v\leq g$。
令 $u_n=s_n(u)$、$v_n=s_n(v)$。它们是简单函数，且
$u_n+v_n\uparrow u+v=s$。引理~\ref{lemma:3-1-1-layer} 的最后一项以及简单函数正线性给出
\[
 \int s
 =\lim_n\int(u_n+v_n)
 =\lim_n\left(\int u_n+\int v_n\right)
 \leq\int f+\int g.
\]
对所有 $s\leq f+g$ 取上确界，便得反向不等式。最后把可加性和齐次性合并即得
\eqref{eq:3-1-positive-linearity}。
\end{proof}

\begin{proposition}[零积分判据]\label{proposition:3-1-1-3}
对非负可测函数 $f$，
\[
 \int f\,\dd\mu=0
 \quad\Longleftrightarrow\quad
 f=0\quad\mu\text{-a.e.}
 \tag{3.1.6}\label{eq:3-1-zero-integral}
\]
因而非负可测函数若 a.e. 相等，其积分相同。
\end{proposition}

\begin{proof}
若 $\int f=0$，令 $E_n=\{f>1/n\}$。因为
$(1/n)\1_{E_n}\leq f$，单调性给出
\[
 0\leq n^{-1}\mu(E_n)\leq\int f=0,
\]
故 $\mu(E_n)=0$。而 $\{f>0\}=\bigcup_nE_n$，所以 $f=0$ a.e.

反过来，设 $f=0$ 于 $X\setminus N$，其中 $\mu(N)=0$。若简单函数 $0\leq s\leq f$，
则 $s$ 的每个正值层都包含于 $N$，故 $\int s=0$。对 $s$ 取上确界即得 $\int f=0$。

最后若 $f=g$ 于 $X\setminus N$，则 $f\1_N$ 与 $g\1_N$ 均由前半部分积分为零；
命题~\ref{proposition:3-1-1-2} 给出
\[
 \int f=\int f\1_{X\setminus N}=\int g\1_{X\setminus N}=\int g.
\]
即使把零函数在一个零测点改成 $+\infty$，这个结论仍然适用。
\end{proof}

现在回算前面的原子模型。对任意 $f\geq0$，规范简单函数满足
$s_n(f)(x)\uparrow f(x)$，故
\[
 \int f\,\dd\delta_x=\sup_ns_n(f)(x)=f(x).
\]
若 $\mu=\sum_{j=1}^Nw_j\delta_{x_j}$ 且各 $x_j$ 互异，则简单函数情形由定义给出
$\int s\,\dd\mu=\sum_jw_js(x_j)$；对 $s_n(f)$ 使用这一等式并令 $n\to\infty$，
有限个非负项可逐项取极限，得到
$\int f\,\dd\mu=\sum_jw_jf(x_j)$。
若 $\mu$ 是计数测度，则简单函数计算和上确界定义给出
\[
 \int f\,\dd\mu
 =\sup\left\{\sum_{x\in F}f(x):F\subset X\text{ 有限}\right\}.
\]
确切地，对有限 $F$ 与 $M>0$，简单函数
$\sum_{x\in F}(f(x)\wedge M)\1_{\{x\}}$ 位于 $f$ 下方；先令 $M\to\infty$，再对 $F$ 取上确界，
给出 ``$\geq$''。
为核对反向不等式，任一 $s\leq f$ 的正值层若有一个无限集，$\int s$ 与右侧上确界都为无穷；
否则 $s$ 的正支撑有限，$\int s=\sum_xs(x)\leq\sup_F\sum_{x\in F}f(x)$。
这就是非负无序和；此处没有任何重排造成的条件抵消。

\begin{theorem}[Radon 开集公式]\label{theorem:3-1-1-4}
设 $X$ 为局部紧 Hausdorff 空间，$\mu$ 为 $X$ 上的 Radon 测度，$U\subset X$ 开。则
\[
 \mu(U)=\sup\left\{\int_X f\,\dd\mu:
 f\in C_c(X),\ 0\leq f\leq1,\ \supp f\subset U\right\}.
 \tag{3.1.7}\label{eq:3-1-radon-open}
\]
\end{theorem}

\begin{proof}
若 $f$ 属于右侧测试类，则 $f\leq\1_U$，所以
$\int f\leq\mu(U)$。反过来，任取紧集 $K\subset U$。局部紧 Hausdorff 空间上的紧支撑截断定理
\ref{theorem:1-3-1-2} 给出 $f\in C_c(X)$，满足
$0\leq f\leq1$、$f=1$ 于 $K$ 且 $\supp f\subset U$。于是
\[
 \mu(K)=\int\1_K\,\dd\mu\leq\int f\,\dd\mu.
\]
Radon 测度在开集上的内正则性给出
$\mu(U)=\sup_{K\subset U,\,K\text{ 紧}}\mu(K)$，故右侧上确界不小于 $\mu(U)$。
\end{proof}

\begin{example}[区间上的帽函数]\label{ex:3-1-1-radon-hat}
取 $X=\R$、$U=(a,b)$，并设 $0<\varepsilon<(b-a)/4$。定义
\[
 f_\varepsilon(x)=
 \begin{cases}
 0,&x\leq a+\varepsilon\ \text{或}\ x\geq b-\varepsilon,\\
 (x-a-\varepsilon)/\varepsilon,&a+\varepsilon<x<a+2\varepsilon,\\
 1,&a+2\varepsilon\leq x\leq b-2\varepsilon,\\
 (b-\varepsilon-x)/\varepsilon,&b-2\varepsilon<x<b-\varepsilon.
 \end{cases}
\]
这是定理~\ref{theorem:3-1-1-4} 中的测试函数。在线性斜坡上用 $N$ 个等长小区间作下阶梯和上阶梯，
相应简单函数的积分分别为
\[
 \frac{\varepsilon}{N}\sum_{j=0}^{N-1}\frac jN
 =\frac{\varepsilon(N-1)}{2N},
 \qquad
 \frac{\varepsilon}{N}\sum_{j=1}^{N}\frac jN
 =\frac{\varepsilon(N+1)}{2N}.
\]
它们分别从下、从上夹住斜坡，故令 $N\to\infty$ 得每条斜坡的积分为 $\varepsilon/2$。因此
\[
 \int f_\varepsilon\,\dd x
 =(b-a-4\varepsilon)+2\cdot\frac{\varepsilon}{2}
 =b-a-3\varepsilon.
\]
为核对所用的单调逼近，令 $d=x-a$。在区间左半部，帽函数可写为
\[
 \max\left\{0,\min\left\{1,\frac d\varepsilon-1\right\}\right\},
\]
该量随 $\varepsilon\downarrow0$ 单调不减，并在 $d>0$ 时趋于 $1$；右半部对
$b-x$ 有同样的表示，而 $x=a,b$ 时始终取零。因此，例如对 $k\geq1$ 取
$\varepsilon_k=(b-a)2^{-k-2}$，就有
\[
 f_{\varepsilon_k}(x)\uparrow\1_{(a,b)}(x)
 \quad(x\in\R),
 \qquad
 \int f_{\varepsilon_k}\,\dd x
 =b-a-3(b-a)2^{-k-2}\longrightarrow b-a.
\]
连续紧支撑测试确实恢复了开区间的长度。
\end{example}

\begin{remark}[拓扑与正则性确实参与了公式]
定理~\ref{theorem:3-1-1-4} 不是任意测度空间上的积分恒等式：没有拓扑时，$C_c(X)$ 根本没有
定义；即使有拓扑，若缺少开集内正则性或把紧集截在开集内的连续帽函数，也无法从
$\mu(K)$ 反向逼近 $\mu(U)$。因此右侧测试类之所以能恢复开集质量，恰好同时使用了
局部紧 Hausdorff 结构与 Radon 正则性。
\end{remark}

\begin{proposition}[非负推前积分公式]\label{proposition:3-1-1-5}
设 $T:(X,\mathcal A)\to(Y,\mathcal B)$ 可测，$\mu$ 为 $X$ 上的测度。对每个非负
$\mathcal B$-可测函数 $f$，
\[
 \int_Y f\,\dd(T_*\mu)=\int_Xf\circ T\,\dd\mu.
 \tag{3.1.8}\label{eq:3-1-pushforward-positive}
\]
\end{proposition}

\begin{proof}
对 $f=\1_B$，两边分别为
$(T_*\mu)(B)$ 与 $\mu(T^{-1}B)$，由推前测度的定义相等。有限正线性继而给出简单函数情形。
对一般 $f$，令 $s_n=s_n(f)$。由于截断、乘法和取整都是逐点操作，
\[
 s_n(f)\circ T=s_n(f\circ T).
\]
引理~\ref{lemma:3-1-1-layer} 分别应用于两侧，再用简单函数情形，得到
\[
 \int_Yf\,\dd(T_*\mu)
 =\sup_n\int_Ys_n(f)\,\dd(T_*\mu)
 =\sup_n\int_Xs_n(f\circ T)\,\dd\mu
 =\int_Xf\circ T\,\dd\mu.
\]
证明并未假定 $f\circ T$ 下方的每个简单函数都来自某个复合函数。
\end{proof}

对经验测度，上式两边都等于 $n^{-1}\sum_{j=1}^nf(Tx_j)$；抽象公式正是有限样本恒等式的延伸。

\begin{figure}[htbp]
\centering
\begin{tikzpicture}[x=1.05cm,y=1.05cm]
  \draw[->,book line] (-0.2,0)--(6.3,0) node[right]{$x$};
  \draw[->,book line] (0,-0.2)--(0,3.7) node[above]{$y$};
  \draw[BookAccent,very thick,smooth] plot coordinates
    {(0,0.2) (0.7,0.65) (1.2,0.55) (1.8,1.45) (2.4,1.2)
     (3.0,2.35) (3.7,2.0) (4.3,3.1) (5.0,2.55) (5.8,3.25)};
  \draw[BookWarm,thick]
    (0,0)--(0.7,0)--(0.7,0.5)--(1.2,0.5)--(1.2,0.5)--(1.8,0.5)
    --(1.8,1.25)--(2.4,1.25)--(2.4,1.0)--(3.0,1.0)--(3.0,2.0)
    --(3.7,2.0)--(3.7,1.75)--(4.3,1.75)--(4.3,2.75)--(5.0,2.75)
    --(5.0,2.5)--(5.8,2.5)--(5.8,0);
  \fill[BookWarm!12] (0,0)--(0.7,0)--(0.7,0.5)--(1.8,0.5)--(1.8,1.25)
    --(2.4,1.25)--(2.4,1)--(3,1)--(3,2)--(3.7,2)--(3.7,1.75)--(4.3,1.75)
    --(4.3,2.75)--(5,2.75)--(5,2.5)--(5.8,2.5)--(5.8,0)--cycle;
  \node[book label,BookAccent] at (4.9,3.45) {$f$};
  \node[book label,BookWarm] at (2.0,0.78) {$s_n(f)$};
\end{tikzpicture}
\caption{非负积分由图像下方的有限阶梯填充；网格越细、截断越高，阶梯单调上升。}
\label{fig:3-1-simple-staircase}
\end{figure}

\begin{remark}[从集合路线到 Daniell 语言]
若一开始给定非负函数锥上的正线性泛函 $I$，自然会令
$\mu(E)=I(\1_E)$，在简单函数上要求
$I(\sum a_k\1_{E_k})=\sum a_k\mu(E_k)$，再以简单函数下逼近的上确界扩张。
上面的构造证明了这套代数骨架在已有测度时是一致的。若要反向推出 Daniell 表示定理，还必须对
$I$ 假设适当的序列连续性，使集合函数 $E\mapsto I(\1_E)$ 具有可数可加性。
\end{remark}

\begin{exercise}
证明若 $s,t$ 是非负简单函数，则 $s\wedge t$ 与 $s\vee t$ 仍为简单函数，并直接从共同分割计算其积分。
\end{exercise}

\begin{exercise}
从两种不交值层表示的公共细分出发，独立重证引理~\ref{lemma:3-1-1-1} 的表示无关性；说明为何
证明中不会出现 $+\infty-(+\infty)$。
\end{exercise}

\begin{exercise}
对 $\mu=\sum_{j=1}^Nw_j\delta_{x_j}$，由定义~\ref{def:3-1-1-2} 证明
$\int f\,\dd\mu=\sum_jw_jf(x_j)$，包括某些 $f(x_j)=+\infty$ 的情形。
\end{exercise}

\begin{exercise}
设 $s_n(f)$ 由 \eqref{eq:3-1-canonical-simple} 定义。证明
$0\leq f-s_n(f)<2^{-n}$ 在 $\{f\leq n\}$ 上成立，并据此给出有界函数的逐点误差界。
\end{exercise}

\begin{exercise}
直接比较相邻两级的截断高度与二进网格，证明
$s_n(f)\le s_{n+1}(f)$，并逐点证明 $s_n(f)\uparrow f$，不得调用 MCT。
\end{exercise}

\begin{exercise}
由定义~\ref{def:3-1-1-2} 重新证明：若 $f\geq0$ 且 $\int f=0$，则
$\mu\{f>\varepsilon\}=0$ 对每个 $\varepsilon>0$，并推出 $f=0$ a.e.
\end{exercise}

从下逼近被写进积分的定义之后，递增极限理应是最容易处理的极限。但 ``理应'' 不是证明；
下一小节将从定义本身建立 MCT，再由它长出 Fatou 与 DCT。

\subsection{MCT、Fatou、DCT 与可数性边界}\label{subsec:3-1-2}
\index{单调收敛定理}\index{Fatou 引理}\index{支配收敛定理}

极限与积分能否交换，是积分理论的第一道压力测试。三种看似相近的函数列给出三种不同答案。
若 $E_n\uparrow E$，指标函数 $\1_{E_n}$ 没有丢失质量；若一个单位质量的区间一路向右移动，
每个固定点最终都看不见它；若质量挤进越来越窄的区间，逐点极限同样可能漏掉全部积分。
因此，逐点收敛不是通行证。单调性、下极限和可积支配分别承担不同的控制任务。
还有一处更隐蔽的断裂：把序列换成任意有向网时，即便各项递增，极限也可能由不可数次拼接产生，
而测度只承诺可数可加。序列定理建立后，我们会给出一个积分始终为零、上确界积分却为一的显式网。

\begin{example}[单调示性函数]\label{ex:3-1-2-1}
若 $E_n\uparrow E$，则 $\1_{E_n}\uparrow\1_E$，而测度从下连续已经告诉我们
$\mu(E_n)\uparrow\mu(E)$。这应当成为一般非负函数的模型。若 $E_n\downarrow E$，
则只有在 $\mu(E_1)<\infty$ 时才能从补集得到 $\mu(E_n)\downarrow\mu(E)$；有限性不是排版上的小字，
而是允许从 $\mu(E_1)$ 中作减法的条件。
\end{example}

\begin{example}[质量逃向无穷远]\label{ex:3-1-2-2}
在 $(\R,m)$ 上令 $f_n=\1_{[n,n+1]}$。对每个 $x$，至多一个 $n$ 使 $f_n(x)=1$，故
$f_n(x)\to0$；但 $\int f_n\,\dd m=1$。于是
\[
 \int\liminf_nf_n\,\dd m=0
 <1=\liminf_n\int f_n\,\dd m.
 \tag{3.1.9}\label{eq:3-1-fatou-strict}
\]
下极限只能保证不凭空增加质量，不能追赶已经跑到无穷远的质量。
\end{example}

\begin{example}[尖峰压缩]\label{ex:3-1-2-3}
在 $(0,1)$ 上令 $f_n=n\1_{(0,1/n)}$。对每个 $x>0$，当 $n>1/x$ 时 $f_n(x)=0$，
所以 $f_n\to0$ 处处；然而
\[
 \int_0^1 f_n\,\dd x=n\,m(0,1/n)=1.
\]
若存在与 $n$ 无关且可积的 $g$ 满足 $f_n\leq g$，下面的支配收敛定理将迫使积分趋零。
所以这个例子不只是说某个证明失效，它证明了统一可积支配函数根本不存在。
\end{example}

我们先证明生成其他结论的根部。

\begin{theorem}[单调收敛定理（MCT）]\label{theorem:3-1-2-1}
设 $f_n:X\to[0,+\infty]$ 可测，并且在一个共同零集之外 $f_n\leq f_{n+1}$。
在该零集上把每个 $f_n$ 改为零，并令 $f$ 为所得序列的逐点上确界。则 $f$ 可测，且
\[
 \int f_n\,\dd\mu\uparrow\int f\,\dd\mu,
 \tag{3.1.10}\label{eq:3-1-mct}
\]
等式允许两边为 $+\infty$；把 $f$ 换成任何与它 a.e. 相等的可测函数，积分不变。
\end{theorem}

\begin{proof}
设 $N$ 是共同异常零集。把每个 $f_n$ 和 $f$ 在 $N$ 上改为零；
命题~\ref{proposition:3-1-1-3} 保证积分不变。此后可以假设 $f_n\uparrow f$ 处处。
由单调性，数列 $\int f_n$ 递增且不超过 $\int f$。记其极限为 $L$。

任取非负简单函数 $s\leq f$ 和 $0<c<1$，令
\[
 E_n=\{x:f_n(x)\geq cs(x)\}.
\]
集合 $E_n$ 递增，并且 $\bigcup_nE_n$ 包含 $\{s>0\}$：若 $s(x)>0$，则
$f_n(x)\uparrow f(x)\geq s(x)>cs(x)$，所以某个 $n$ 后 $x\in E_n$。
由 $f_n\geq cs\1_{E_n}$ 与积分单调性，
\[
 \int f_n\,\dd\mu\geq c\int s\1_{E_n}\,\dd\mu.
\]
把 $s$ 写成有限个正值层 $\sum_{j=1}^ra_j\1_{A_j}$。对每个 $j$，
$A_j\cap E_n\uparrow A_j$；测度从下连续给出
\[
 \lim_n\int s\1_{E_n}\,\dd\mu
 =\sum_{j=1}^ra_j\lim_n\mu(A_j\cap E_n)
 =\int s\,\dd\mu.
\]
因此 $L\geq c\int s$。令 $c\uparrow1$，再对全部 $s\leq f$ 取上确界，得到
$L\geq\int f$。结合先前的反向不等式即得 \eqref{eq:3-1-mct}。

证明中唯一把无穷多个阶段合并起来的地方，是对递增的集合列使用测度从下连续；
这也标出了可数性进入证明的位置。
\end{proof}

\begin{corollary}[非负级数可逐项积分]\label{corollary:3-1-2-series}
若 $u_k:X\to[0,+\infty]$ 可测，则
\[
 \int_X\sum_{k=1}^{\infty}u_k\,\dd\mu
 =\sum_{k=1}^{\infty}\int_Xu_k\,\dd\mu,
 \tag{3.1.10a}\label{eq:3-1-series-integral}
\]
两边允许为 $+\infty$。
\end{corollary}

\begin{proof}
部分和 $S_n=\sum_{k=1}^n u_k$ 递增至 $S=\sum_{k=1}^\infty u_k$。
MCT 与非负积分的有限可加性依次给出
\[
 \int S=\lim_n\int S_n
 =\lim_n\sum_{k=1}^n\int u_k
 =\sum_{k=1}^\infty\int u_k.
\]
\end{proof}

\begin{corollary}[集合的单调连续性]\label{corollary:3-1-2-set-continuity}
若 $E_n\uparrow E$，则 $\mu(E_n)\uparrow\mu(E)$。若 $E_n\downarrow E$ 且
$\mu(E_1)<\infty$，则 $\mu(E_n)\downarrow\mu(E)$。
\end{corollary}

\begin{proof}
第一项对 $\1_{E_n}\uparrow\1_E$ 应用 MCT。第二项令
$F_n=E_1\setminus E_n$，则 $F_n\uparrow E_1\setminus E$，故
\[
 \mu(E_1)-\mu(E_n)=\mu(F_n)\uparrow\mu(E_1\setminus E)=\mu(E_1)-\mu(E).
\]
由于 $\mu(E_1)<\infty$，这些减法都发生在有限实数中。
\end{proof}

有限性不能删去。在 $\R$ 上取 $E_n=[n,+\infty)$，则 $E_n\downarrow\varnothing$，
但 $m(E_n)=+\infty$ 对所有 $n$ 成立，因而这些测度不可能下降到 $m(\varnothing)=0$。

回到例~\ref{ex:3-1-2-1}，示性函数计算现在已成为一般定理的直接特例。
下一条结论不要求函数列单调；代价是它只保留下界。

\begin{theorem}[Fatou 引理]\label{theorem:3-1-2-2}
若 $f_n:X\to[0,+\infty]$ 可测，则
\[
 \int_X\liminf_{n\to\infty}f_n\,\dd\mu
 \leq\liminf_{n\to\infty}\int_Xf_n\,\dd\mu.
 \tag{3.1.11}\label{eq:3-1-fatou}
\]
\end{theorem}

\begin{proof}
令 $g_n=\inf_{k\geq n}f_k$。这是可数族可测函数的下确界，因而可测，并且
$g_n\uparrow\liminf_kf_k$。MCT 给出
\[
 \int\liminf_kf_k=\lim_n\int g_n.
\]
对每个 $k\geq n$ 有 $g_n\leq f_k$，故
\[
 \int g_n\leq\inf_{k\geq n}\int f_k.
\]
令 $n\to\infty$ 即得 \eqref{eq:3-1-fatou}。
\end{proof}

例~\ref{ex:3-1-2-2} 说明 Fatou 不等式可以严格。若另有一个可积上界，
把同一不等式同时用于 ``剩余质量''，便能夹出等号。

\begin{theorem}[非负支配收敛定理]\label{theorem:3-1-2-3}
设 $g:X\to[0,\infty]$ 为可测函数，$0\leq f_n\leq g$ a.e.，且
$f_n\to f$ a.e.，其中 $f$ 可测并满足 $\int g\,\dd\mu<\infty$。则
$0\leq f\leq g$ a.e.，并且
\[
 \int|f_n-f|\,\dd\mu\longrightarrow0,
 \qquad
 \int f_n\,\dd\mu\longrightarrow\int f\,\dd\mu.
 \tag{3.1.12}\label{eq:3-1-nonnegative-dct}
\]
\end{theorem}

\begin{proof}
先令 $Z=\{g=+\infty\}$。对每个 $M>0$，有 $M\1_Z\le g$，因而
$M\mu(Z)\le\int g<\infty$；令 $M\to\infty$ 得 $\mu(Z)=0$。把 $Z$ 与收敛、支配失败的
可数多个零集并成一个零集，再在其上把 $g,f_n,f$ 全部改为零。积分不变，此后 $g$ 处处有限，
$0\leq f_n\leq g$ 且 $f_n\to f$ 处处，故 $0\leq f\leq g$。
特别地，$\int f_n$ 与 $\int f$ 都不超过有限数 $\int g$。

Fatou 引理先给出
\[
 \int f\leq\liminf_n\int f_n.
 \tag{3.1.13}\label{eq:3-1-dct-liminf}
\]
再对非负函数 $g-f_n\to g-f$ 应用 Fatou。非负积分的可加性以及 $\int g<\infty$ 给出
\[
 \int g-\int f
 \leq\liminf_n\left(\int g-\int f_n\right)
 =\int g-\limsup_n\int f_n.
\]
在有限实数中消去 $\int g$，得到
$\limsup_n\int f_n\leq\int f$。与 \eqref{eq:3-1-dct-liminf} 合并，便有积分收敛。

还需证明绝对误差趋零。函数
$h_n=2g-|f_n-f|$ 非负且处处趋于 $2g$。Fatou 引理给出
\[
 2\int g
 \leq\liminf_n\int h_n
 =\liminf_n\left(2\int g-\int|f_n-f|\right)
 =2\int g-\limsup_n\int|f_n-f|.
\]
再次在有限实数中消去 $2\int g$，得到
$\limsup_n\int|f_n-f|\leq0$。这些积分非负，所以它们趋于零。
\end{proof}

例~\ref{ex:3-1-2-3} 的尖峰由此排除了任何统一可积支配函数。
请注意，本定理只积分非负函数；允许负值和复值的版本要等到下一小节定义相应积分后才能陈述。

\begin{example}[一次参数积分的回收计算]\label{ex:3-1-2-parameter}
令
\[
 F(t)=\int_0^1e^{-tx}\,\dd x\qquad(t\in\R).
\]
若 $t_n\to t$，则存在 $M$ 使 $|t_n|\leq M$，从而
$0<e^{-t_nx}\leq e^M$ 于 $[0,1]$，而常数 $e^M$ 在该区间可积。
非负 DCT 给出 $F(t_n)\to F(t)$。由于 $\R$ 是第一可数的，序列连续性等价于连续性，故 $F$ 连续；
一元微积分的直接计算还给出
\[
 F(t)=
 \begin{cases}
 (1-e^{-t})/t,&t\ne0,\\
 1,&t=0.
 \end{cases}
\]
这里极限定理把逐点连续提升成了积分后的连续性；同一机制也适用于由核给出的参数积分。
\end{example}

\begin{definition}[一致尾控制]\label{def:3-1-2-3}
设 $\mathcal F$ 是一族实值或复值可测函数，且每个 $f\in\mathcal F$ 都满足
$\int|f|<\infty$。若
\[
 \lim_{K\to\infty}\sup_{f\in\mathcal F}
 \int_{\{|f|>K\}}|f|\,\dd\mu=0,
 \tag{3.1.14}\label{eq:3-1-uniform-tail}
\]
则称 $\mathcal F$ 具有一致尾控制。
\end{definition}

式~\eqref{eq:3-1-uniform-tail} 只涉及非负函数 $|f|$ 的积分，因而定义已经完备。
尖峰族 $f_n=n\1_{(0,1/n)}$ 不满足该条件：取 $n>K$ 时，
$\int_{|f_n|>K}|f_n|=1$。不过，在无限测度空间上，一致尾控制连一致 $L^1$ 有界性也不能单独保证。
例如 $f_n=\1_{[0,n]}$ 在 $K>1$ 时使 \eqref{eq:3-1-uniform-tail} 的上确界为零，
但 $\int f_n=n$。即使另加一致 $L^1$ 界，高度尾控制仍不阻止空间逃逸：
$g_n=\1_{[n,n+1]}$ 满足 $\|g_n\|_1=1$，且在 $K>1$ 时尾积分恒为零，但对每个紧集
$C\subset\R$ 都有 $\int_{\R\setminus C}g_n=1$ 对所有充分大的 $n$ 成立。
因此一致尾控制只限制函数的高度尾部；在无限测度空间上，还必须分别检查
$L^1$ 有界性和防止质量逃向无穷远的紧性条件。

序列极限定理背后是测度的可数可加性。拓扑中常用网描述极限，那么 MCT 能否把 $n$ 换成任意有向指标？
下面的例子把答案钉死。

\begin{example}[不可数网陷阱]\label{ex:3-1-2-4}
设 $X$ 是不可数集合，$\mathcal A$ 是可数--余可数 $\sigma$-代数，并定义
\[
 \mu(A)=
 \begin{cases}
 0,&A\text{ 可数},\\
 1,&X\setminus A\text{ 可数}.
 \end{cases}
\]
这是概率测度。以 $X$ 的所有有限子集组成的集合 $I$ 为有向集，次序取包含关系，并令
$f_F=\1_F$。则 $F\subset G$ 蕴含 $f_F\leq f_G$，而逐点上
$\sup_{F\in I}f_F=\1_X$。然而每个 $F$ 有 $\int f_F=\mu(F)=0$，而
$\int\1_X=1$。因此任意网版本的 MCT 是假的。
\end{example}

这个有向集没有可数共尾子集：若 $(F_n)$ 共尾，则可数集 $\bigcup_nF_n$ 应包含每个单点，
从而等于不可数的 $X$，矛盾。安全的推广恰好止于能够降回序列的网。

不可数方向还可能更早出故障。取一个非 Borel 集 $A\subset\R$，并令有限子集
$F\subset A$ 按包含关系成网。每个 $\1_F$ 都是 Borel 可测的，但
\[
 \sup_{F\subset A,\,F\text{ 有限}}\1_F=\1_A
\]
不是 Borel 可测函数。可数个可测函数的上确界可测，并不蕴含任意族也有这一性质。

\begin{definition}[可数共尾网]\label{def:3-1-2-2}
有向集 $I$ 若含有可数共尾子集，即存在可数 $C\subset I$，使对每个 $i\in I$ 都有
$c\in C$ 满足 $i\leq c$，则称 $I$ 可数共尾。以这样的 $I$ 为指标的网简称可数共尾网。
\end{definition}

第一可数性与这里的概念不同。第一可数性说某个拓扑空间在每一点有可数邻域基；
可数共尾性说索引集内部有一条可数的最终路径。值域即使是 $\R$，
例~\ref{ex:3-1-2-4} 的索引集仍然没有可数共尾子集。

\begin{proposition}[可数共尾网版本]\label{proposition:3-1-2-4}
设 $I$ 为可数共尾有向集。
\begin{enumerate}
 \item 若 $(f_i)_{i\in I}$ 是递增的非负可测网，且 $f=\sup_{i\in I}f_i$ 可测，则
 \[
   \sup_{i\in I}\int f_i\,\dd\mu=\int f\,\dd\mu.
 \]
 \item 若 $0\leq f_i\leq g$ a.e.，$\int g<\infty$，且 $f_i\to f$ a.e.，其中 $f$ 可测，
 则 $\int|f_i-f|\to0$。
\end{enumerate}
这里的 a.e. 陈述均使用一个对全网共同有效的零集。一般网不在本命题中享有无条件的 Fatou 引理。
\end{proposition}

\begin{proof}
枚举一个可数共尾集为 $(c_n)$。利用有向性递归选取 $d_n$，使
$d_1\geq c_1$ 且 $d_n\geq d_{n-1},c_n$。于是 $(d_n)$ 是递增共尾序列。
第一项中，$f_{d_n}\uparrow f$：它们的上确界不超过 $f$；反过来，对任意 $i$，
共尾性给出某个 $c_n\geq i$，继而 $d_n\geq c_n$，所以 $f_i\leq f_{d_n}$。
序列 MCT 因而给出
\[
 \int f=\lim_n\int f_{d_n}=\sup_{i\in I}\int f_i.
\]

对第二项，若结论失败，则存在 $\varepsilon>0$，使对每个 $i_0\in I$ 都能找到
$i\geq i_0$ 满足 $\int|f_i-f|\geq\varepsilon$。先选 $i_1\ge d_1$ 保持这一坏性质；已选 $i_{n-1}$ 后，
再递归选取 $i_n\geq d_n,i_{n-1}$ 且保持坏性质。于是 $(i_n)$ 递增，且因 $i_n\ge d_n$ 而共尾；
映射 $n\mapsto i_n$ 保序且共尾，所以它按子网的定义确实给出原网的一个子网。因此
$f_{i_n}\to f$ a.e.，且仍由 $g$ 支配。序列的非负支配收敛定理给出
$\int|f_{i_n}-f|\to0$，与其始终不小于 $\varepsilon$ 矛盾。

从可数共尾集构造递增链以及逐步选取坏指标，是可数次递归选择；按照第一章的选择约定，
这些步骤使用可数依赖选择。一般网的尾下确界可能是不可数族的下确界，甚至未必可测，
所以没有额外的可数尾决定条件时，不能把 Fatou 的序列证明照抄过去。
\end{proof}

\begin{figure}[htbp]
\centering
\begin{tikzpicture}[node distance=10mm and 13mm]
  \node[book strong node] (mct) {MCT\\递增极限};
  \node[book node,above left=of mct] (sets) {集合从下连续\\从上连续需有限性};
  \node[book node,above right=of mct] (fatou) {Fatou\\尾下确界};
  \node[book warm node,right=of fatou] (dct) {非负 DCT\\可积支配};
  \node[book warning node,below right=of mct] (net) {任意网断裂\\不可数拼接};
  \draw[book accent arrow] (mct)--(sets);
  \draw[book accent arrow] (mct)--(fatou);
  \draw[book accent arrow] (fatou)--(dct);
  \draw[book dashed arrow] (mct)--(net);
  \node[book label,below=2mm of dct] {$2g-|f_n-f|$ 控制误差};
\end{tikzpicture}
\caption{极限定理的生成关系。虚线分支不是另一个定理，而是不可数网反例标出的边界。}
\label{fig:3-1-convergence-tree}
\end{figure}

\begin{remark}[三个直接用途]
推论~\ref{corollary:3-1-2-series} 已经给出级数与积分交换的非负版本，
例~\ref{ex:3-1-2-parameter} 则展示 DCT 如何产生参数积分连续性。
第三种用途是把 Cauchy 列抽成绝对增量可求和的快速子列：若
$\sum_k\|f_{n_{k+1}}-f_{n_k}\|_p<\infty$，则
\eqref{eq:3-1-series-integral} 能把范数估计转化为几乎处处的绝对收敛。这正是
Riesz--Fischer 完备性证明的核心机制。
\end{remark}

\begin{exercise}
利用推论~\ref{corollary:3-1-2-series} 计算
$\int_0^1\sum_{k=0}^\infty x^k\,\dd x$，并解释为什么答案为 $+\infty$。
\end{exercise}

\begin{exercise}
设 $0\leq f_n\leq g$ 且 $\int g<\infty$。用 Fatou 引理分别作用于 $f_n$ 与 $g-f_n$，
证明
\[
 \limsup_n\int f_n\leq\int\limsup_nf_n.
\]
说明这里为什么需要逐点不等式 $f_n\leq g$。
\end{exercise}

\begin{exercise}
构造三个函数列，分别说明非负 DCT 中 ``a.e. 收敛''、``同一个支配函数'' 和
``支配函数可积'' 不能任意删去。
\end{exercise}

\begin{exercise}
完整证明例~\ref{ex:3-1-2-4} 中的 $\mu$ 是可数可加测度，并证明有限子集有向集没有可数共尾子集。
\end{exercise}

非负极限定理已经允许我们稳定地处理大小，却仍不允许抵消。
$+\infty-(+\infty)$ 不是一个等待聪明记号来修补的空格；下一小节将精确规定何时减法合法。

\subsection{实值与复值积分：抵消、绝对值与局部化}\label{subsec:3-1-3}
\index{正部与负部}\index{$L^1$ 空间@$L^1$ 空间}\index{局部可积}

在 $\R$ 上看函数 $f(x)=x$。每个对称截断 $[-R,R]$ 都给出零，但在整个实线上，
正半轴和负半轴各自携带无穷质量。若把这两个无穷硬写成 ``相减为零''，换一种截断方式就会得到别的答案。
对称主值是一种附加的求和规则，不是 Lebesgue 积分。我们需要的原则是：先分别测量正负质量，
只有在减法有意义时才允许抵消。

\begin{definition}[正负部分]\label{def:3-1-3-1}
对扩展实值可测函数 $f$，定义
\[
 f^+=\max\{f,0\},\qquad f^-=\max\{-f,0\}.
\]
于是 $f=f^+-f^-$ 且 $|f|=f^++f^-$。若 $\int f^+$ 与 $\int f^-$ 不同时为
$+\infty$，定义
\[
 \int f\,\dd\mu:=\int f^+\,\dd\mu-\int f^-\,\dd\mu.
 \tag{3.1.15}\label{eq:3-1-signed-integral}
\]
若两者都有限，则称 $f$ 可积。
\end{definition}

\begin{definition}[$L^1$ 可积]\label{def:3-1-3-2}
实值或复值可测函数 $f$ 若满足
\[
 \int_X|f|\,\dd\mu<\infty,
\]
就称为可积。所有这样的点值函数组成 $\mathcal L^1(\mu)$；按 a.e. 相等取商所得空间记为
$L^1(\mu)$。对复值 $f$ 定义
\[
 \int f\,\dd\mu:=\int\Re f\,\dd\mu+i\int\Im f\,\dd\mu,
 \qquad
 \|[f]\|_1:=\int|f|\,\dd\mu.
 \tag{3.1.16}\label{eq:3-1-complex-integral}
\]
不致混淆时省略等价类的方括号。
\end{definition}

因为 $|\Re f|,|\Im f|\leq|f|$，复积分右侧确实由两个有限实数构成。
命题~\ref{proposition:3-1-1-3} 保证 $\|[f]\|_1=0$ 当且仅当 $[f]=0$；
所以它在等价类空间上才是范数，在点值函数空间上只是半范数。

\begin{example}[条件抵消不可定义]\label{ex:3-1-3-1}
对 $f(x)=x$，把 $[0,R]$ 等分成 $N$ 段。以左、右端点高度构成的简单函数分别夹住 $x$，其积分为
\[
 \frac{R^2(N-1)}{2N}
 \leq\int_0^R x\,\dd x
 \leq\frac{R^2(N+1)}{2N}.
\]
令 $N\to\infty$ 得 $\int_0^Rx\,\dd x=R^2/2$。
因此 $\int_\R f^+=\int_0^\infty x\,\dd x=+\infty$；同样
$\int_\R f^-=+\infty$。故 $\int_\R x\,\dd x$ 没有定义。
另一方面，每个 $[-R,R]$ 上的积分为零：把两侧分别用等长阶梯计算，正负部分的有限积分相等。
这说明 ``所有对称截断都为零'' 并不能替代定义~\ref{def:3-1-3-1}。
\end{example}

\begin{definition}[局部可积]\label{def:3-1-3-3}
设 $X$ 为局部紧 Hausdorff 空间，$\mu$ 为 Radon 测度。若可测函数 $f$ 对每个紧集
$K\subset X$ 都满足
\[
 \int_K|f|\,\dd\mu<\infty,
\]
则称 $f$ 局部可积，记作 $f\in L^1_{\mathrm{loc}}(X,\mu)$。
\end{definition}

若 $X$ 还第二可数，第一章给出的紧耗尽满足
$K_n\subset\operatorname{int}K_{n+1}$ 且 $\bigcup_nK_n=X$。此时只须检查每个 $K_n$：
任意紧集 $K$ 被开覆盖 $\{\operatorname{int}K_n\}$ 覆盖，有限子覆盖与递增性给出
$K\subset K_N$。

\begin{example}[局部可积但不可积]\label{ex:3-1-3-2}
相对于 $\R^n$ 上的 Lebesgue 测度，常数函数 $1$ 在每个紧集 $K$ 上满足
$\int_K1=m_n(K)<\infty$，故属于 $L^1_{\mathrm{loc}}(\R^n)$。但
\[
 \int_{\R^n}1\,\dd x=m_n(\R^n)=+\infty,
\]
所以它不属于 $L^1(\R^n)$。局部化允许我们在每个有限窗口内工作，并没有把无穷空间伪装成有限空间。
\end{example}

\begin{example}[幂奇点的局部可积阈值]\label{ex:3-1-3-power}
在 $\R^n$ 上令 $f(x)=|x|^{-\alpha}$，其中 $\alpha\geq0$，并令 $f(0)=+\infty$。
只需检查原点附近。置
\[
 A_k=\{x:2^{-k-1}<|x|\leq2^{-k}\},\qquad k\geq0.
\]
若 $v_n=m_n(B(0,1))$，第~2~章的线性体积缩放给出
\[
 m_n(A_k)=v_n(1-2^{-n})2^{-kn}.
\]
在 $A_k$ 上有
$2^{k\alpha}\leq f(x)<2^{(k+1)\alpha}$。对不交环带的有限并使用非负可加性，
再令环带数目增加并用 MCT，得到
\[
 C_1\sum_{k=0}^N2^{-k(n-\alpha)}
 \leq\int_{\bigcup_{k=0}^NA_k}|x|^{-\alpha}\,\dd x
 \leq C_2\sum_{k=0}^N2^{-k(n-\alpha)},
 \tag{3.1.17}\label{eq:3-1-power-annuli}
\]
其中 $C_1=v_n(1-2^{-n})$、$C_2=2^\alpha C_1$。
几何级数收敛当且仅当 $\alpha<n$。原点是零测单点，远离原点时函数在紧集上有界，故
\[
 |x|^{-\alpha}\in L^1_{\mathrm{loc}}(\R^n)
 \quad\Longleftrightarrow\quad \alpha<n.
\]
极坐标公式也会给出同一阈值；上面的环带论证则只使用体积缩放与几何级数。
\end{example}

\begin{proposition}[限制测度与集合上积分]\label{proposition:3-1-3-restriction}
对 $E\in\mathcal A$ 定义限制测度
$(\mu|_E)(A)=\mu(A\cap E)$。若 $f$ 非负可测或可积，则
\[
 \int_Ef\,\dd\mu
 =\int_Xf\1_E\,\dd\mu
 =\int_Xf\,\dd(\mu|_E).
 \tag{3.1.18}\label{eq:3-1-restriction}
\]
\end{proposition}

\begin{proof}
第一式是记号的定义。若 $f=\1_A$，则后两式都是 $\mu(A\cap E)$；有限正线性给出简单函数情形。
对一般非负 $f$，$s_n(f)$ 在两种测度下都是规范简单逼近，故由
引理~\ref{lemma:3-1-1-layer}，
\[
 \int f\,\dd(\mu|_E)
 =\sup_n\int s_n(f)\,\dd(\mu|_E)
 =\sup_n\int s_n(f)\1_E\,\dd\mu
 =\int f\1_E\,\dd\mu.
\]
实值可积函数拆成正负部分，复值函数拆成实虚部，即得一般情形。
\end{proof}

经验测度下，\eqref{eq:3-1-restriction} 的三式都等于
$n^{-1}\sum_{j:x_j\in E}f(x_j)$，限制测度就是把不符合筛选条件的样本权重置零。

下面建立 $L^1$ 的代数。证明不能先假定 ``积分差由差的积分控制''，因为这句话本身需要积分线性。

\begin{theorem}[$L^1$ 上积分的线性]\label{theorem:3-1-3-2}
$L^1(\mu)$ 是实或复向量空间，并且
\[
 I:L^1(\mu)\to\R\ \text{或}\ \C,
 \qquad I(f)=\int f\,\dd\mu,
\]
是线性泛函。a.e. 相等的可积函数有相同积分。
\end{theorem}

\begin{proof}
若 $f,g\in\mathcal L^1$，则点态三角不等式和非负积分的可加性给出
\[
 \int|f+g|\leq\int(|f|+|g|)=\int|f|+\int|g|<\infty;
\]
对任意标量 $a$，$\int|af|=|a|\int|f|<\infty$。因此可积函数形成向量空间，a.e. 等价与
向量运算相容。

先设 $f,g$ 实值。点态恒等式
\[
 (f+g)^+ +f^-+g^-
 =(f+g)^-+f^++g^+
 \tag{3.1.19}\label{eq:3-1-linearity-identity}
\]
两边都是非负可积函数。对两边积分并使用命题~\ref{proposition:3-1-1-2}，
所有项均为有限实数，故可以移项，得到
\[
 \int(f+g)=\int f+\int g.
\]
当 $a\geq0$ 时，$(af)^+=af^+$、$(af)^-=af^-$；当 $a<0$ 时，
正负部分交换并乘以 $|a|$。所以 $\int af=a\int f$。

对复值函数，实虚部分别应用上述结论，先得加法性。又
\[
 \Re(if)=-\Im f,
 \qquad \Im(if)=\Re f,
\]
故 $\int(if)=i\int f$；结合实齐次性便得复线性。

若 $f=g$ a.e.，则 $|f-g|=0$ a.e.，命题~\ref{proposition:3-1-1-3} 给出
$\int|f-g|=0$。在实值情形，$0\le(f-g)^\pm\le|f-g|$，故两个正负部分的积分都为零，
从定义得到 $\int(f-g)=0$；复值情形对实部、虚部重复此论证，因为二者的绝对值都不超过
$|f-g|$。由已经在点值可积函数上证明的线性，$\int f-\int g=\int(f-g)=0$。因此积分确实
下降为 $L^1$ 等价类上的线性泛函。
\end{proof}

\begin{example}[积分平均与中心化]\label{ex:3-1-3-expectation}
设 $(\Omega,\mathcal F,\Prob)$ 是总质量为一的测度空间。对可积可测函数 $X$，记其积分平均为
$\operatorname{Av}_{\Prob}(X):=\int_\Omega X\,\dd\Prob$；在概率论中，这一数值称为期望。常数函数
$c$ 的积分是
$c\Prob(\Omega)=c$，故积分线性给出
\[
 \operatorname{Av}_{\Prob}\!\left(X-\operatorname{Av}_{\Prob}(X)\right)
 =\operatorname{Av}_{\Prob}(X)
  -\operatorname{Av}_{\Prob}(X)\Prob(\Omega)=0.
\]
因此 ``中心化'' 不是形式上的减号，而是由总质量为一和 $L^1$ 线性共同保证的零平均操作。
\end{example}

\begin{proposition}[绝对值估计]\label{proposition:3-1-3-1}
对每个 $f\in L^1(\mu)$，
\[
 \left|\int f\,\dd\mu\right|\leq\int|f|\,\dd\mu.
 \tag{3.1.20}\label{eq:3-1-integral-absolute}
\]
\end{proposition}

\begin{proof}
若 $f$ 实值，则 $-|f|\leq f\leq|f|$。积分线性与非负积分单调性给出
$-\int|f|\leq\int f\leq\int|f|$。

若 $f$ 复值且 $z=\int f\ne0$，取单位复数 $\lambda=\overline z/|z|$，使
$\lambda z=|z|$。由已经证明的复线性和实值情形，
\[
 \left|\int f\right|
 =\Re\left(\lambda\int f\right)
 =\int\Re(\lambda f)
 \leq\int|\lambda f|
 =\int|f|.
\]
当 $z=0$ 时结论成立。
\end{proof}

\begin{example}[复相位只改变抵消]\label{ex:3-1-3-3}
设 $g$ 可测，$\theta:X\to\R$ 可测，并令 $f=e^{i\theta}g$。连续复合保证
$e^{i\theta}$ 可测，而且 $|f|=|g|$，所以
\[
 \int|f|=\int|g|
\]
允许共同为 $+\infty$。若这一数值有限，则 $f,g\in L^1$，但一般
$\int f\ne\int g$：相位可以制造抵消，却不会改变控制量。命题~\ref{proposition:3-1-3-1}
说明绝对值恰是所有相位共同服从的上界。
\end{example}

\begin{corollary}[实值与复值推前积分公式]\label{corollary:3-1-3-pushforward}
设 $T:X\to Y$ 可测。若 $f:Y\to\R$ 或 $\C$ 对 $T_*\mu$ 可积，则
$f\circ T$ 对 $\mu$ 可积，且
\[
 \int_Yf\,\dd(T_*\mu)=\int_Xf\circ T\,\dd\mu.
 \tag{3.1.21}\label{eq:3-1-pushforward-signed}
\]
\end{corollary}

\begin{proof}
对非负函数 $|f|$ 使用命题~\ref{proposition:3-1-1-5}，得到
\[
 \int_X|f\circ T|\,\dd\mu
 =\int_Y|f|\,\dd(T_*\mu)<\infty.
\]
故复合函数可积。实值时分别对 $f^+,f^-$ 使用非负推前公式；复值时再对实部、虚部使用实值结论。
\end{proof}

\begin{theorem}[实值与复值支配收敛]\label{theorem:3-1-3-5}
设 $f_n:X\to\R$ 或 $\C$ 可测，$f_n\to f$ a.e.，其中 $f$ 可测。若存在非负可测函数
$g$ 满足 $|f_n|\leq g$ a.e. 且 $\int g<\infty$，则 $f,f_n\in L^1(\mu)$，并且
\[
 \|f_n-f\|_1\longrightarrow0,
 \qquad
 \int f_n\,\dd\mu\longrightarrow\int f\,\dd\mu.
 \tag{3.1.22}\label{eq:3-1-full-dct}
\]
\end{theorem}

\begin{proof}
共同化异常零集并在其上把所有函数改为零后，可设收敛与支配处处成立。
取极限得 $|f|\leq g$，故 $f,f_n$ 可积；并且
$|f_n-f|\leq2g$ 且 $|f_n-f|\to0$。对非负函数列 $|f_n-f|$ 应用
定理~\ref{theorem:3-1-2-3}，得到 $\|f_n-f\|_1\to0$。最后由积分线性和绝对值估计，
\[
 \left|\int f_n-\int f\right|
 =\left|\int(f_n-f)\right|
 \leq\int|f_n-f|\longrightarrow0.
\]
\end{proof}

这就是实值或复值可积函数的支配收敛定理（DCT）。

\begin{proposition}[限制与拼接]\label{proposition:3-1-3-3}
若 $(E_k)_{k\geq1}$ 是两两不交的可测集且 $f\in L^1(\mu)$，则
\[
 \int_{\bigcup_{k=1}^\infty E_k}f\,\dd\mu
 =\sum_{k=1}^\infty\int_{E_k}f\,\dd\mu,
 \tag{3.1.23}\label{eq:3-1-patching}
\]
右侧级数绝对收敛。
\end{proposition}

\begin{proof}
对非负函数 $|f|$，有限部分和
$\sum_{k=1}^n|f|\1_{E_k}$ 递增至 $|f|\1_{\cup_kE_k}$。MCT 与有限可加性给出
\[
 \sum_{k=1}^\infty\int_{E_k}|f|
 =\int_{\cup_kE_k}|f|
 \leq\int|f|<\infty.
\]
由命题~\ref{proposition:3-1-3-1}，
$|\int_{E_k}f|\leq\int_{E_k}|f|$，故右侧绝对收敛。
对每个 $n$，积分线性给出
\[
 \sum_{k=1}^n\int_{E_k}f=\int_{\cup_{k=1}^nE_k}f.
\]
两边与目标值之差的绝对值不超过
$\int_{\cup_{k>n}E_k}|f|$；该量由上面的收敛级数趋于零。令 $n\to\infty$ 即得公式。
\end{proof}

有限复测度的积分还要解决另一个问题：积分变量本身携带相位。第~2~章由分割定义的全变差
$|\nu|$ 能直接控制简单函数积分，标量 $L^1(|\nu|)$ 的完备性继而给出延拓。

\begin{theorem}[复测度积分估计]\label{theorem:3-1-3-4}
设 $\nu$ 是 $(X,\mathcal A)$ 上的有限复测度。若
$s=\sum_{k=1}^ma_k\1_{E_k}$ 是复值简单函数，其中 $E_k$ 两两不交，定义
\[
 \int_Xs\,\dd\nu:=\sum_{k=1}^ma_k\nu(E_k).
\]
这个定义与表示无关，并且可唯一延伸到每个 $f\in L^1(|\nu|)$，使
\[
 \left|\int_Xf\,\dd\nu\right|
 \leq\int_X|f|\,\dd|\nu|.
 \tag{3.1.24}\label{eq:3-1-complex-measure-bound}
\]
延伸后的积分在 $L^1(|\nu|)$ 上复线性。
\end{theorem}

\begin{proof}
表示无关的证明与引理~\ref{lemma:3-1-1-1} 相同：取公共有限分割，并使用复测度的有限可加性。
由全变差的定义和 $|\nu(E)|\leq|\nu|(E)$，
\[
 \left|\int s\,\dd\nu\right|
 \leq\sum_k|a_k|\,|\nu(E_k)|
 \leq\sum_k|a_k|\,|\nu|(E_k)
 =\int|s|\,\dd|\nu|.
 \tag{3.1.25}\label{eq:3-1-simple-complex-measure-bound}
\]

给定 $f\in L^1(|\nu|)$，分别以规范层级逼近 $\Re f$、$\Im f$ 的正负部分：
\[
 s_n=\bigl(s_n((\Re f)^+)-s_n((\Re f)^-)\bigr)
   +i\bigl(s_n((\Im f)^+)-s_n((\Im f)^-)\bigr).
\]
这些是复值简单函数。MCT 给出
$\int s_n(h)\uparrow\int h$ 对每个非负可积的 $h$；由非负可加性，
\[
 \int(h-s_n(h))\,\dd|\nu|
 =\int h\,\dd|\nu|-\int s_n(h)\,\dd|\nu|\longrightarrow0.
\]
四项相加并用点态三角不等式，得到
\[
 \int|s_n-f|\,\dd|\nu|\longrightarrow0.
 \tag{3.1.26}\label{eq:3-1-simple-l1-approx}
\]
由 \eqref{eq:3-1-simple-complex-measure-bound}，
\[
 \left|\int s_n\,\dd\nu-\int s_m\,\dd\nu\right|
 \leq\int|s_n-s_m|\,\dd|\nu|,
\]
所以 $\bigl(\int s_n\,\dd\nu\bigr)$ 是 $\C$ 中的 Cauchy 列。利用 $\C$ 的完备性，定义
\[
 \int f\,\dd\nu:=\lim_n\int s_n\,\dd\nu.
\]
若 $(t_n)$ 是另一列满足 $\int|t_n-f|\,\dd|\nu|\to0$ 的简单函数，
则
\[
 \left|\int s_n\,\dd\nu-\int t_n\,\dd\nu\right|
 \leq\int|s_n-f|\,\dd|\nu|+\int|t_n-f|\,\dd|\nu|\to0,
\]
故极限与逼近选择无关。对简单逼近逐项使用线性并取极限，得到复线性。

最后，$\bigl||s_n|-|f|\bigr|\leq|s_n-f|$，所以
$\int|s_n|\,\dd|\nu|\to\int|f|\,\dd|\nu|$。在
\eqref{eq:3-1-simple-complex-measure-bound} 中令 $n\to\infty$，便得
\eqref{eq:3-1-complex-measure-bound}。若另一个延伸在线性且受同一 $L^1$ 误差控制，
它在简单函数上相同，并由 \eqref{eq:3-1-simple-l1-approx} 在每个 $f$ 上相同，故延伸唯一。
整个构造只使用了标量域的完备性，没有调用 Hahn--Banach、$L^p$ 完备性或复测度的极分解。
\end{proof}

\begin{example}[两个复原子]\label{ex:3-1-3-complex-atoms}
令
\[
 \nu=i\delta_0-(1+i)\delta_1.
\]
分割定义给出 $|\nu|=\delta_0+\sqrt2\,\delta_1$。因此对任意函数 $f:\{0,1\}\to\C$，
\[
 \int f\,\dd\nu=if(0)-(1+i)f(1),
\]
且
\[
 \left|if(0)-(1+i)f(1)\right|
 \leq |f(0)|+\sqrt2|f(1)|
 =\int|f|\,\dd|\nu|.
\]
这里相位与控制量都能逐原子看见。
\end{example}

\begin{figure}[htbp]
\centering
\begin{tikzpicture}[x=1cm,y=1cm]
  \draw[->,book line] (-3.8,0)--(3.8,0) node[right]{$x$};
  \draw[->,book line] (0,-2.4)--(0,2.7) node[above]{$y$};
  \draw[BookAccent,very thick,smooth] plot coordinates
    {(-3,0.2) (-2.4,1.25) (-1.7,1.75) (-1.0,0.55) (-0.3,0)
     (0.35,-0.8) (1.1,-1.65) (1.8,-0.7) (2.5,0) (3.1,1.0)};
  \fill[BookAccent!12] (-3,0)--(-3,0.2)--(-2.4,1.25)--(-1.7,1.75)--(-1,0.55)--(-0.3,0)--cycle;
  \fill[BookWarning!12] (0.35,0)--(0.35,-0.8)--(1.1,-1.65)--(1.8,-0.7)--(2.5,0)--cycle;
  \draw[BookWarning,dashed,thick,smooth] plot coordinates
    {(0.35,0.8) (1.1,1.65) (1.8,0.7) (2.5,0)};
  \node[book label,BookAccent] at (-1.8,2.15) {$f^+$};
  \node[book label,BookWarning] at (1.2,-2.0) {$f^-$};
  \node[book label,BookWarning] at (1.3,2.1) {$|f|$ 将负区翻上};
\end{tikzpicture}
\caption{带符号积分是正区质量减去负区质量；绝对值把两区都变成非负控制量。}
\label{fig:3-1-positive-negative-parts}
\end{figure}

\begin{remark}[复测度的极表示]
定理~\ref{theorem:3-1-3-4} 直接从全变差定义了 $\int f\,\dd\nu$。稍后的
定理~\ref{theorem:3-3-1-3} 还将给出 $\nu=u|\nu|$，其中 $|u|=1$ a.e.；由此可把同一积分写成
\[
  \int f\,\dd\nu=\int fu\,\dd|\nu|.
\]
前一种定义强调全变差控制，后一种表示则把复测度积分化为正测度积分。
\end{remark}

\begin{remark}[等价类没有点态支撑]
$L^1$ 元素是 a.e. 等价类。在零测稠密集上修改一个代表，可能彻底改变其点态非零集的闭包，
却不改变 $L^1$ 元素、范数或积分。因此，若要谈函数的拓扑支撑，必须指定代表；
若只谈 $L^1$ 结论，则应使用关于零集修改稳定的本质支撑概念。
\end{remark}

\begin{exercise}
证明 $\|f\|_1=\int|f|$ 在 $L^1(\mu)$ 上满足三角不等式，并说明为什么在
$\mathcal L^1(\mu)$ 上可能只有半范数。
\end{exercise}

\begin{exercise}
分别给出 $\R$ 上一个属于 $L^1_{\mathrm{loc}}\setminus L^1$ 的有界函数和一个无界函数，
并按定义验证。
\end{exercise}

\begin{exercise}
若 $f\in L^1$ 且 $E_n\downarrow\varnothing$，证明
$\int_{E_n}|f|\to0$。提示：对 $|f|\1_{E_n}$ 使用 DCT。
\end{exercise}

\begin{exercise}
证明定理~\ref{theorem:3-1-3-4} 中简单函数积分的表示无关；再证明若
$f_n\to f$ 于 $L^1(|\nu|)$，则 $\int f_n\,\dd\nu\to\int f\,\dd\nu$。
\end{exercise}

我们现在拥有两套互补规则：Tonelli 型问题应先保持非负，Fubini 型问题则要用绝对可积性控制抵消。
下一节把它们放到积空间中，并检验两种积分次序何时能够交换。
